\documentclass[12pt]{article}
\usepackage[margin=1in]{geometry}
\usepackage{amsmath}
\usepackage{booktabs}
\usepackage{graphicx}

\title{K-Means Image Segmentation: Implementation Report}
\author{Nathaniel Black \\ CISC484: Introduction to Machine Learning \\ Professor Xu Yuan}
\date{April 14, 2026}

\begin{document}
\maketitle

\section{Introduction}

Before writing any code, I started by reading Stuart Lloyd's 1982 paper
``Least Squares Quantization in PCM.'' Although Lloyd originally wrote this as
a Bell Labs internal memo in 1957 for scalar quantization of audio signals,
the core algorithm generalizes directly to higher-dimensional clustering and is
the foundation of what we now call k-means. I focused my reading on Sections
IV, V, and VI, which contain the full algorithm, and skimmed Appendix B, which
demonstrates why initialization matters.

Section IV proves that the optimal representative for a fixed group of points
is the center of mass (the mean), which corresponds to the centroid update
step. Section V proves that the optimal partition given fixed representatives
is to assign each point to the nearest representative, which corresponds to
the assignment step. Section VI combines these into an alternating iterative
procedure (Method I) and proves that the total squared error decreases
monotonically at each iteration, guaranteeing convergence to a local minimum.

The connection between Lloyd's 1D quantization problem and image segmentation
is that both are instances of the same optimization: given a set of points
(pixel colors in RGB space, or voltage levels on a line), find $k$
representative values and assign each point to one of them such that the total
squared error is minimized. In the image case, each pixel is a 3D vector
$(R, G, B)$, the $k$ centroids are the representative colors, and the
``quantized'' image replaces each pixel with its assigned centroid color.

\section{Task 1: Scikit-learn Implementation}

The scikit-learn implementation uses the \texttt{KMeans} class directly. The
image is loaded with PIL, converted to a numpy array of shape $(H, W, 3)$, and
reshaped to $(N, 3)$ where $N = H \times W$ so that each row is one pixel's
RGB values. \texttt{KMeans} is called with \texttt{n\_clusters=k},
\texttt{random\_state=42} for reproducibility, and \texttt{n\_init=10} to run
the algorithm 10 times with different initializations and keep the best
result. The segmented image is reconstructed by indexing into
\texttt{cluster\_centers\_} with the predicted labels and reshaping back to
$(H, W, 3)$.

\section{Task 2: K-Means from Scratch}

I built two from-scratch implementations: a non-optimized version using Python
objects and loops, and an optimized version using vectorized numpy operations.

\subsection{Non-optimized Version}

This version closely mirrors the algorithm as described in Lloyd's paper, step
by step, using explicit Python data structures. I defined a \texttt{Pixel}
class with \texttt{r}, \texttt{g}, \texttt{b} attributes and a
\texttt{centroid} attribute tracking which cluster the pixel belongs to. I
implemented \texttt{\_\_hash\_\_} and \texttt{\_\_eq\_\_} on the
\texttt{Pixel} class so that \texttt{Pixel} objects could be used as
dictionary keys, which was necessary for grouping pixels by their assigned
centroid using a \texttt{defaultdict}.

The algorithm proceeds as follows:
\begin{enumerate}
    \item All pixels in the image are converted to \texttt{Pixel} objects.
    \item $k$ random \texttt{Pixel} objects are selected as initial centroids
          using uniform random sampling.
    \item In each iteration, every pixel computes its Euclidean distance to all
          centroids and is assigned to the nearest one (the assignment step from
          Section V).
    \item Pixels are grouped by their centroid using a \texttt{defaultdict},
          and the mean R, G, B values of each group become the new centroids
          (the update step from Section IV).
    \item Steps 3--4 repeat until the set of centroids stops changing between
          iterations.
    \item The segmented image is reconstructed by replacing each pixel with its
          centroid's color.
\end{enumerate}

\subsection{Optimized Version}

This version performs the same algorithm but uses numpy broadcasting to
vectorize the computation. Instead of \texttt{Pixel} objects, pixels remain as
a numpy array of shape $(N, 3)$ cast to \texttt{float64} to avoid
\texttt{uint8} overflow during subtraction. Initial centroids are selected
using \texttt{np.random.choice} to pick $k$ unique random indices.

The key optimization is in the distance computation. Instead of looping over
pixels and centroids, the expression
\begin{verbatim}
np.linalg.norm(pixels[:, np.newaxis] - centroids[np.newaxis, :], axis=2)
\end{verbatim}
computes all $N \times k$ distances in a single vectorized operation by
broadcasting the $(N, 1, 3)$ pixel array against the $(1, k, 3)$ centroid
array. \texttt{np.argmin} along axis 1 then gives each pixel's nearest
centroid index.

The centroid update is also vectorized: for each cluster $i$,
\texttt{pixels[labels == i].mean(axis=0)} computes the mean of all pixels in
that cluster. If a cluster becomes empty (no pixels assigned to it), the old
centroid is retained to avoid \texttt{NaN} values from computing the mean of
an empty slice.

Convergence is checked using \texttt{np.array\_equal} on the centroid arrays.
When centroids stop changing between iterations, the loop breaks.

\section{Comparison of Results}

The \texttt{main.py} script runs all three implementations (scikit-learn,
non-optimized, and optimized) on all four provided images (beacon, elephant,
robin, vulture) and displays the results in a grid for side-by-side comparison.

The scikit-learn implementation produces the most consistent results because it
uses \texttt{n\_init=10}, running the algorithm 10 times and keeping the result
with the lowest total squared error. This mitigates the initialization
sensitivity that Lloyd noted in Appendix B of his paper.

The optimized numpy version produces results visually comparable to
scikit-learn. The non-optimized version tends to show slightly more color
variance because it uses integer floor division (\texttt{//}) when computing
centroid means, which introduces rounding bias that accumulates over
iterations. The float-based implementations (scikit-learn and the optimized
version) do not have this issue.

All three implementations produce valid color quantizations of the input
images, reducing them to $k$ distinct colors while preserving the overall
structure and recognizable features of the original images.

\subsection{Timing Results}

I added timing measurements to compare the performance of all three
implementations across each image with $k = 5$.

\begin{table}[h]
\centering
\begin{tabular}{lccc}
\toprule
Image & Scikit-learn & Non-optimized & Optimized \\
\midrule
beacon.JPEG   & 0.786s & 0.245s & 0.562s \\
elephant.jpg  & 0.713s & 0.222s & 0.487s \\
robin.JPEG    & 0.764s & 0.256s & 1.357s \\
vulture.JPEG  & 0.750s & 0.245s & 0.868s \\
\bottomrule
\end{tabular}
\caption{Execution time for each implementation with $k = 5$.}
\end{table}

The results were surprising. The non-optimized Python loop version was
consistently the fastest, while scikit-learn was among the slowest.
Scikit-learn's slower time is straightforward to explain: \texttt{n\_init=10}
means it runs the full algorithm 10 times and keeps the best result, so it is
doing roughly 10 times the work of a single run. The optimized numpy version
being slower than the Python loop version is less obvious. For these
relatively small images, numpy's broadcasting allocates large intermediate
arrays of shape $(N, k, 3)$ for the distance computation, and the memory
allocation overhead outweighs the speed gains from vectorized arithmetic. On
larger images where $N$ is much bigger, the vectorized version would be
expected to pull ahead as the computation-to-overhead ratio shifts in its
favor.

\section{Bugs Encountered During Development}

Several bugs came up during implementation that are worth noting because they
reflect real pitfalls of translating an algorithm from paper to code:

\begin{itemize}
    \item \textbf{Initial centroid sampling range:} I initially sampled random
          indices from \texttt{range(0, k)} instead of
          \texttt{range(0, len(pixels))}, meaning I was only ever selecting
          from the first few pixels in the image. This resulted in all
          centroids being nearly identical and the output being a single flat
          color.
    \item \textbf{Missing closest-distance update:} In the nearest centroid
          function, I checked \texttt{if distance < closest} but forgot to
          update \texttt{closest = distance} inside the if block. This caused
          every pixel to be assigned to the last centroid evaluated rather than
          the nearest one.
    \item \textbf{uint8 overflow:} The optimized version initially operated on
          \texttt{uint8} pixel values. Subtracting two \texttt{uint8} values
          wraps around (e.g., $5 - 10$ becomes $251$ instead of $-5$),
          producing incorrect distances and visible banding artifacts in the
          output. Casting to \texttt{float64} fixed this.
    \item \textbf{Empty clusters:} With some initializations, certain clusters
          end up with no pixels assigned. Computing the mean of an empty set
          produces \texttt{NaN}, which propagates through subsequent
          iterations. Retaining the old centroid for empty clusters resolves
          this.
    \item \textbf{Convergence check:} I initially set
          \texttt{centroids = new\_centroids} before checking for convergence,
          making the old and new centroids the same object and causing the loop
          to terminate after one iteration. The fix was to compare before
          overwriting.
\end{itemize}

\section{Lessons Learned}

Reading Lloyd's original paper before implementing was valuable. The paper
makes clear that k-means is solving a least-squares optimization problem
through coordinate descent: the assignment step optimizes the partition with
centroids held fixed, and the update step optimizes the centroids with the
partition held fixed. Each step can only decrease the objective, which is why
the algorithm converges. Understanding this made debugging easier because I
could check whether my total squared error was monotonically decreasing.

The timing results challenged my assumption that vectorized numpy code is
always faster than Python loops. For small images, the memory overhead of
allocating large intermediate arrays can dominate the computation time,
making the simpler loop-based implementation faster in practice. This
reinforced that optimization decisions depend on the actual problem size, not
just the asymptotic complexity of the approach.

\end{document}
